\documentclass[10pt,a4paper]{amsart}
\usepackage{amsmath,amssymb,amsthm,graphicx}
% Retained minimal preamble for the Stacks Project chapter sites-cohomology.tex (Cohomology on Sites).
% Verbatim excerpt of the chapter preamble: Xy-pic setup (native preamble lines 14-18) and the native
% theorem declarations with the mathematical macros (native preamble lines 154-217). Copied mechanically;
% nothing added, removed or redefined. The wrapper class and the amsmath/amssymb/amsthm/graphicx base
% packages are supplied by the assembler, so the class/externaldocument part of the native preamble is omitted.
% For commutative diagrams we use Xy-pic
\usepackage[all]{xy}

% We use 2cell for 2-commutative diagrams.
\xyoption{2cell}
% Theorem environments.
%
\theoremstyle{plain}
\newtheorem{theorem}[subsection]{Theorem}
\newtheorem{proposition}[subsection]{Proposition}
\newtheorem{lemma}[subsection]{Lemma}

\theoremstyle{definition}
\newtheorem{definition}[subsection]{Definition}
\newtheorem{example}[subsection]{Example}
\newtheorem{exercise}[subsection]{Exercise}
\newtheorem{situation}[subsection]{Situation}

\theoremstyle{remark}
\newtheorem{remark}[subsection]{Remark}
\newtheorem{remarks}[subsection]{Remarks}

\numberwithin{equation}{subsection}

% Macros
%
\def\lim{\mathop{\mathrm{lim}}\nolimits}
\def\colim{\mathop{\mathrm{colim}}\nolimits}
\def\Spec{\mathop{\mathrm{Spec}}}
\def\Hom{\mathop{\mathrm{Hom}}\nolimits}
\def\Ext{\mathop{\mathrm{Ext}}\nolimits}
\def\SheafHom{\mathop{\mathcal{H}\!\mathit{om}}\nolimits}
\def\SheafExt{\mathop{\mathcal{E}\!\mathit{xt}}\nolimits}
\def\Sch{\mathit{Sch}}
\def\Mor{\mathop{\mathrm{Mor}}\nolimits}
\def\Ob{\mathop{\mathrm{Ob}}\nolimits}
\def\Sh{\mathop{\mathit{Sh}}\nolimits}
\def\NL{\mathop{N\!L}\nolimits}
\def\CH{\mathop{\mathrm{CH}}\nolimits}
\def\proetale{{pro\text{-}\acute{e}tale}}
\def\etale{{\acute{e}tale}}
\def\QCoh{\mathit{QCoh}}
\def\Ker{\mathop{\mathrm{Ker}}}
\def\Im{\mathop{\mathrm{Im}}}
\def\Coker{\mathop{\mathrm{Coker}}}
\def\Coim{\mathop{\mathrm{Coim}}}

% Boxtimes
%
\DeclareMathSymbol{\boxtimes}{\mathbin}{AMSa}{"02}

%
% Macros for moduli stacks/spaces
%
\def\QCohstack{\mathcal{QC}\!\mathit{oh}}
\def\Cohstack{\mathcal{C}\!\mathit{oh}}
\def\Spacesstack{\mathcal{S}\!\mathit{paces}}
\def\Quotfunctor{\mathrm{Quot}}
\def\Hilbfunctor{\mathrm{Hilb}}
\def\Curvesstack{\mathcal{C}\!\mathit{urves}}
\def\Polarizedstack{\mathcal{P}\!\mathit{olarized}}
\def\Complexesstack{\mathcal{C}\!\mathit{omplexes}}
% \Pic is the operator that assigns to X its picard group, usage \Pic(X)
% \Picardstack_{X/B} denotes the Picard stack of X over B
% \Picardfunctor_{X/B} denotes the Picard functor of X over B
\def\Pic{\mathop{\mathrm{Pic}}\nolimits}
\def\Picardstack{\mathcal{P}\!\mathit{ic}}
\def\Picardfunctor{\mathrm{Pic}}
\def\Deformationcategory{\mathcal{D}\!\mathit{ef}}
\usepackage{cleveref}
\crefname{originallemma}{original lemma}{original lemmas}
\Crefname{originallemma}{Original lemma}{Original lemmas}
\makeatletter
\newlabel{lemma-cohomology-bigger-site}{{03YU}{source}{}{sourcepointer-1}{}}
\newlabel{lemma-cohomology-bigger-site@cref}{{[originallemma][0][]03YU}{[1][0][]source}{}{}{}}
\makeatother
\crefname{originallemma}{original lemma}{original lemmas}
\Crefname{originallemma}{Original lemma}{Original lemmas}
\makeatletter
\newlabel{sites-lemma-bigger-site}{{00XU}{source}{}{sourcepointer-2}{}}
\newlabel{sites-lemma-bigger-site@cref}{{[originallemma][0][]00XU}{[1][0][]source}{}{}{}}
\makeatother
\crefname{originallemma}{original lemma}{original lemmas}
\Crefname{originallemma}{Original lemma}{Original lemmas}
\makeatletter
\newlabel{lemma-describe-pullback-pi}{{09X3}{source}{}{sourcepointer-3}{}}
\newlabel{lemma-describe-pullback-pi@cref}{{[originallemma][0][]09X3}{[1][0][]source}{}{}{}}
\makeatother
\crefname{originallemma}{original lemma}{original lemmas}
\Crefname{originallemma}{Original lemma}{Original lemmas}
\makeatletter
\newlabel{sites-lemma-relocalize}{{03EH}{source}{}{sourcepointer-4}{}}
\newlabel{sites-lemma-relocalize@cref}{{[originallemma][0][]03EH}{[1][0][]source}{}{}{}}
\makeatother
\crefname{originallemma}{original lemma}{original lemmas}
\Crefname{originallemma}{Original lemma}{Original lemmas}
\makeatletter
\newlabel{sites-lemma-relocalize-given-fibre-products}{{09W9}{source}{}{sourcepointer-5}{}}
\newlabel{sites-lemma-relocalize-given-fibre-products@cref}{{[originallemma][0][]09W9}{[1][0][]source}{}{}{}}
\makeatother
\crefname{originallemma}{original lemma}{original lemmas}
\Crefname{originallemma}{Original lemma}{Original lemmas}
\makeatletter
\newlabel{sites-lemma-composition-morphisms-sites}{{03CB}{source}{}{sourcepointer-6}{}}
\newlabel{sites-lemma-composition-morphisms-sites@cref}{{[originallemma][0][]03CB}{[1][0][]source}{}{}{}}
\makeatother
\crefname{originallemma}{original lemma}{original lemmas}
\Crefname{originallemma}{Original lemma}{Original lemmas}
\makeatletter
\newlabel{lemma-cohomology-of-open}{{03F3}{source}{}{sourcepointer-7}{}}
\newlabel{lemma-cohomology-of-open@cref}{{[originallemma][0][]03F3}{[1][0][]source}{}{}{}}
\makeatother
\crefname{originallemma}{original lemma}{original lemmas}
\Crefname{originallemma}{Original lemma}{Original lemmas}
\makeatletter
\newlabel{cohomology-lemma-proper-base-change-sheaves-of-sets}{{0D90}{source}{}{sourcepointer-8}{}}
\newlabel{cohomology-lemma-proper-base-change-sheaves-of-sets@cref}{{[originallemma][0][]0D90}{[1][0][]source}{}{}{}}
\makeatother
\crefname{originallemma}{original lemma}{original lemmas}
\Crefname{originallemma}{Original lemma}{Original lemmas}
\makeatletter
\newlabel{lemma-unbounded-describe-higher-direct-images}{{0D6I}{source}{}{sourcepointer-9}{}}
\newlabel{lemma-unbounded-describe-higher-direct-images@cref}{{[originallemma][0][]0D6I}{[1][0][]source}{}{}{}}
\makeatother
\crefname{originallemma}{original lemma}{original lemmas}
\Crefname{originallemma}{Original lemma}{Original lemmas}
\makeatletter
\newlabel{cohomology-lemma-before-Leray}{{01EZ}{source}{}{sourcepointer-10}{}}
\newlabel{cohomology-lemma-before-Leray@cref}{{[originallemma][0][]01EZ}{[1][0][]source}{}{}{}}
\makeatother
\crefname{originallemma}{original lemma}{original lemmas}
\Crefname{originallemma}{Original lemma}{Original lemmas}
\makeatletter
\newlabel{cohomology-theorem-proper-base-change}{{09V6}{source}{}{sourcepointer-11}{}}
\newlabel{cohomology-theorem-proper-base-change@cref}{{[originallemma][0][]09V6}{[1][0][]source}{}{}{}}
\makeatother
\title{The comparison between the Zariski topoi of the big and small sites of a locally quasi-compact space}
\author{The Stacks Project Authors}
\date{stacks-project:sites-cohomology.tex:local-snapshot-20261011; GNU FDL 1.2 or later}
\begin{document}
\maketitle
\noindent\textit{Editorial note.} All delivered mathematical text is the author's own, copied from the pinned Stacks Project chapter file sites-cohomology.tex; no statement or proof step was rewritten, reordered or removed.\par\medskip
\noindent\textit{Editorial note.} The proof is the author's complete argument for the seven items, with two Xy-pic diagrams delivered as written (the square of topoi in the statement and the fibred square in the proof of (7a)).\par\medskip
\noindent\textit{Editorial note.} The label occurs in no other staged Stacks chapter; the checked list is chow, cohomology, divisors, etale-cohomology, flat, more-algebra, more-morphisms, morphisms, spaces-more-morphisms, varieties together with this card's chapter sites-cohomology.\par\medskip
\noindent\textit{Editorial note.} Every labelled reference occurring inside the counted spans points to a chapter of the Project that is not part of this card's source tree. Each is carried in assembly.external\_references as a non-load-bearing navigation pointer whose source location is the invocation inside the counted span, so no load-bearing mathematics is replaced by a stub.\par\medskip
\noindent\textit{Editorial note.} The preamble is the chapter's own retained minimal preamble (sources/stacks-sites-cohomology/retained-preamble.tex), a verbatim excerpt of the native preamble: the Xy-pic setup and the native theorem declarations with the mathematical macros.\par\medskip
\noindent\textit{Allowed standard prerequisites (editorial).}\begin{itemize}
\item The category LC\_Zar of locally quasi-compact Hausdorff spaces with the Zariski topology and the morphisms of topoi pi\_X. Same chapter, the Zariski topology on LC, the topos of sheaves on LC\_Zar/X and the morphism pi\_X relating it to the usual topos of sheaves on X. Supplies every object and morphism occurring in the statement.
\item Cohomology of the big and small Zariski sites compared by an adjunction. Sites, the results comparing a site with a larger site along a continuous and cocontinuous functor, applied here to the inclusion of the open subsets of X into LC\_Zar/X. Supplies the identifications pi\_\{X,*\} = tau\_X\textasciicircum{}\{-1\} and the identity tau\_X\textasciicircum{}\{-1\} pi\_X\textasciicircum{}\{-1\} = id from which parts (1)-(3) are read off.
\item Higher direct images on the big Zariski site and their description by a presheaf. Same chapter, the lemma describing the higher direct images R\textasciicircum{}q f\_\{Zar,*\} by the sheaf associated to a presheaf of cohomology groups on the slices. Supplies the computation in part (7b), where the cohomology sheaves of Rf\_\{Zar,*\}pi\_X\textasciicircum{}\{-1\}K are identified through that description.
\end{itemize}\medskip
\ifdefined\hypertarget\hypertarget{sourcepointer-1}{}\fi
\noindent\textit{Source reference.} Original lemma 03YU: sources/stacks-sites-cohomology/pinned-original-source.tex, source lines 646--646. Navigation pointer invoked at proof line 7812; the counted span contains the pointer only and the cited result is not reproduced inside it, so the pointer is registered as a non-load-bearing navigation stub in the established form for Stacks material. No bibliography citation is involved.\par\medskip
\ifdefined\hypertarget\hypertarget{sourcepointer-2}{}\fi
\noindent\textit{Source reference.} Original lemma 00XU: sources/stacks-sites-cohomology/pinned-original-source.tex, source lines 7818--7818. Navigation pointer invoked at proof line 7818; the counted span contains the pointer only and the cited result is not reproduced inside it, so the pointer is registered as a non-load-bearing navigation stub in the established form for Stacks material. No bibliography citation is involved.\par\medskip
\ifdefined\hypertarget\hypertarget{sourcepointer-3}{}\fi
\noindent\textit{Source reference.} Original lemma 09X3: sources/stacks-sites-cohomology/pinned-original-source.tex, source lines 7675--7675. Navigation pointer invoked at proof line 7826; the counted span contains the pointer only and the cited result is not reproduced inside it, so the pointer is registered as a non-load-bearing navigation stub in the established form for Stacks material. No bibliography citation is involved.\par\medskip
\ifdefined\hypertarget\hypertarget{sourcepointer-4}{}\fi
\noindent\textit{Source reference.} Original lemma 03EH: sources/stacks-sites-cohomology/pinned-original-source.tex, source lines 7833--7833. Navigation pointer invoked at proof line 7833; the counted span contains the pointer only and the cited result is not reproduced inside it, so the pointer is registered as a non-load-bearing navigation stub in the established form for Stacks material. No bibliography citation is involved.\par\medskip
\ifdefined\hypertarget\hypertarget{sourcepointer-5}{}\fi
\noindent\textit{Source reference.} Original lemma 09W9: sources/stacks-sites-cohomology/pinned-original-source.tex, source lines 7836--7836. Navigation pointer invoked at proof line 7836; the counted span contains the pointer only and the cited result is not reproduced inside it, so the pointer is registered as a non-load-bearing navigation stub in the established form for Stacks material. No bibliography citation is involved.\par\medskip
\ifdefined\hypertarget\hypertarget{sourcepointer-6}{}\fi
\noindent\textit{Source reference.} Original lemma 03CB: sources/stacks-sites-cohomology/pinned-original-source.tex, source lines 7839--7839. Navigation pointer invoked at proof line 7839; the counted span contains the pointer only and the cited result is not reproduced inside it, so the pointer is registered as a non-load-bearing navigation stub in the established form for Stacks material. No bibliography citation is involved.\par\medskip
\ifdefined\hypertarget\hypertarget{sourcepointer-7}{}\fi
\noindent\textit{Source reference.} Original lemma 03F3: sources/stacks-sites-cohomology/pinned-original-source.tex, source lines 608--608. Navigation pointer invoked at proof line 7846; the counted span contains the pointer only and the cited result is not reproduced inside it, so the pointer is registered as a non-load-bearing navigation stub in the established form for Stacks material. No bibliography citation is involved.\par\medskip
\ifdefined\hypertarget\hypertarget{sourcepointer-8}{}\fi
\noindent\textit{Source reference.} Original lemma 0D90: sources/stacks-sites-cohomology/pinned-original-source.tex, source lines 7877--7877. Navigation pointer invoked at proof line 7877; the counted span contains the pointer only and the cited result is not reproduced inside it, so the pointer is registered as a non-load-bearing navigation stub in the established form for Stacks material. No bibliography citation is involved.\par\medskip
\ifdefined\hypertarget\hypertarget{sourcepointer-9}{}\fi
\noindent\textit{Source reference.} Original lemma 0D6I: sources/stacks-sites-cohomology/pinned-original-source.tex, source lines 4434--4434. Navigation pointer invoked at proof line 7887; the counted span contains the pointer only and the cited result is not reproduced inside it, so the pointer is registered as a non-load-bearing navigation stub in the established form for Stacks material. No bibliography citation is involved.\par\medskip
\ifdefined\hypertarget\hypertarget{sourcepointer-10}{}\fi
\noindent\textit{Source reference.} Original lemma 01EZ: sources/stacks-sites-cohomology/pinned-original-source.tex, source lines 7907--7907. Navigation pointer invoked at proof line 7907; the counted span contains the pointer only and the cited result is not reproduced inside it, so the pointer is registered as a non-load-bearing navigation stub in the established form for Stacks material. No bibliography citation is involved.\par\medskip
\ifdefined\hypertarget\hypertarget{sourcepointer-11}{}\fi
\noindent\textit{Source reference.} Original lemma 09V6: sources/stacks-sites-cohomology/pinned-original-source.tex, source lines 7909--7909. Navigation pointer invoked at proof line 7909; the counted span contains the pointer only and the cited result is not reproduced inside it, so the pointer is registered as a non-load-bearing navigation stub in the established form for Stacks material. No bibliography citation is involved.\par\medskip
\begin{lemma}
\label{lemma-collect-true-things-Zar}
Let $X$ be an object of $\textit{LC}_{Zar}$. Then
\begin{enumerate}
\item for $\mathcal{F} \in \textit{Ab}(X)$ we have
$H^n_{Zar}(X, \pi_X^{-1}\mathcal{F}) = H^n(X, \mathcal{F})$,
\item $\pi_{X, *} : \textit{Ab}(\textit{LC}_{Zar}/X) \to \textit{Ab}(X)$
is exact,
\item the unit $\text{id} \to \pi_{X, *} \circ \pi_X^{-1}$
of the adjunction is an isomorphism, and
\item for $K \in D(X)$ the canonical map
$K \to R\pi_{X, *} \pi_X^{-1}K$ is an isomorphism.
\end{enumerate}
Let $f : X \to Y$ be a morphism of $\textit{LC}_{Zar}$. Then
\begin{enumerate}
\item[(5)] there is a commutative diagram
$$
\xymatrix{
\Sh(\textit{LC}_{Zar}/X) \ar[r]_{f_{Zar}} \ar[d]_{\pi_X} &
\Sh(\textit{LC}_{Zar}/Y) \ar[d]^{\pi_Y} \\
\Sh(X_{Zar}) \ar[r]^f &
\Sh(Y_{Zar})
}
$$
of topoi,
\item[(6)] for $L \in D^+(Y)$ we have
$H^n_{Zar}(X, \pi_Y^{-1}L) = H^n(X, f^{-1}L)$,
\item[(7)] if $f$ is proper, then we have
\begin{enumerate}
\item $\pi_Y^{-1} \circ f_* = f_{Zar, *} \circ \pi_X^{-1}$ as functors
$\Sh(X) \to \Sh(\textit{LC}_{Zar}/Y)$,
\item $\pi_Y^{-1} \circ Rf_* = Rf_{Zar, *} \circ \pi_X^{-1}$ as
functors $D^+(X) \to D^+(\textit{LC}_{Zar}/Y)$.
\end{enumerate}
\end{enumerate}
\end{lemma}

\begin{proof}
Proof of (1).
The equality $H^n_{Zar}(X, \pi_X^{-1}\mathcal{F}) = H^n(X, \mathcal{F})$
is a general fact coming from the trivial observation that
coverings of $X$ in $\textit{LC}_{Zar}$ are the same thing as open
coverings of $X$. The reader who wishes to see a detailed proof
should apply Lemma \ref{lemma-cohomology-bigger-site} to the functor
$X_{Zar} \to \textit{LC}_{Zar}$.

\medskip\noindent
Proof of (2). This is true because $\pi_{X, *} = \tau_X^{-1}$
for some morphism of topoi $\tau_X : \Sh(X_{Zar}) \to \Sh(\textit{LC}_{Zar})$
as follows from Sites, Lemma \ref{sites-lemma-bigger-site}
applied to the functor
$X_{Zar} \to \textit{LC}_{Zar}/X$ used to define $\pi_X$.

\medskip\noindent
Proof of (3). This is true because $\tau_X^{-1} \circ \pi_X^{-1}$
is the identity functor by Sites, Lemma \ref{sites-lemma-bigger-site}.
Or you can deduce it from the explicit description of
$\pi_X^{-1}$ in Lemma \ref{lemma-describe-pullback-pi}.

\medskip\noindent
Proof of (4). Apply (3) to an complex of abelian sheaves representing $K$.

\medskip\noindent
Proof of (5). The morphism of topoi $f_{Zar}$ comes from an application of
Sites, Lemma \ref{sites-lemma-relocalize}
and in our case comes from the continuous functor
$Z/Y \mapsto Z \times_Y X/X$ by
Sites, Lemma \ref{sites-lemma-relocalize-given-fibre-products}.
The diagram commutes simply because the corresponding
continuous functors compose correctly
(see Sites, Lemma \ref{sites-lemma-composition-morphisms-sites}).

\medskip\noindent
Proof of (6). We have
$H^n_{Zar}(X, \pi_Y^{-1}\mathcal{G}) =
H^n_{Zar}(X, f_{Zar}^{-1}\pi_Y^{-1}\mathcal{G})$
for $\mathcal{G}$ in $\textit{Ab}(Y)$, see
Lemma \ref{lemma-cohomology-of-open}.
This is equal to $H^n_{Zar}(X, \pi_X^{-1}f^{-1}\mathcal{G})$
by the commutativity of the diagram in (5).
Hence we conclude by (1) in the case $L$ consists of a single
sheaf in degree $0$. The general case follows by representing
$L$ by a bounded below complex of abelian sheaves.

\medskip\noindent
Proof of (7a). Let $\mathcal{F}$ be a sheaf on $X$.
Let $g : Z \to Y$ be an object of $\textit{LC}_{Zar}/Y$. Consider the
fibre product
$$
\xymatrix{
Z' \ar[r]_{f'} \ar[d]_{g'} & Z \ar[d]^g \\
X \ar[r]^f & Y
}
$$
Then we have
$$
(f_{Zar, *}\pi_X^{-1}\mathcal{F})(Z/Y) =
(\pi_X^{-1}\mathcal{F})(Z'/X)  =
\Gamma(Z', (g')^{-1}\mathcal{F})  =
\Gamma(Z, f'_*(g')^{-1}\mathcal{F})
$$
the second equality by Lemma \ref{lemma-describe-pullback-pi}.
On the other hand
$$
(\pi_Y^{-1}f_*\mathcal{F})(Z/Y) = \Gamma(Z, g^{-1}f_*\mathcal{F})
$$
again by Lemma \ref{lemma-describe-pullback-pi}.
Hence by proper base change for sheaves of sets
(Cohomology, Lemma \ref{cohomology-lemma-proper-base-change-sheaves-of-sets})
we conclude the two sets are canonically isomorphic.
The isomorphism is compatible with restriction mappings
and defines an isomorphism
$\pi_Y^{-1}f_*\mathcal{F} = f_{Zar, *}\pi_X^{-1}\mathcal{F}$.
Thus an isomorphism of functors
$\pi_Y^{-1} \circ f_* = f_{Zar, *} \circ \pi_X^{-1}$.

\medskip\noindent
Proof of (7b). Let $K \in D^+(X)$. By
Lemma \ref{lemma-unbounded-describe-higher-direct-images}
the $n$th cohomology sheaf of
$Rf_{Zar, *}\pi_X^{-1}K$ is the sheaf associated to the presheaf
$$
(g : Z \to Y) \longmapsto H^n_{Zar}(Z', \pi_X^{-1}K)
$$
with notation as above. Observe that
\begin{align*}
H^n_{Zar}(Z', \pi_X^{-1}K)
& =
H^n(Z', (g')^{-1}K) \\
& =
H^n(Z, Rf'_*(g')^{-1}K) \\
& =
H^n(Z, g^{-1}Rf_*K) \\
& =
H^n_{Zar}(Z, \pi_Y^{-1}Rf_*K)
\end{align*}
The first equality is (6) applied to $K$ and $g' : Z' \to X$.
The second equality is Leray for $f' : Z' \to Z$
(Cohomology, Lemma \ref{cohomology-lemma-before-Leray}).
The third equality is the proper base change theorem
(Cohomology, Theorem \ref{cohomology-theorem-proper-base-change}).
The fourth equality is (6) applied to $g : Z \to Y$ and $Rf_*K$.
Thus $Rf_{Zar, *}\pi_X^{-1}K$ and $\pi_Y^{-1}Rf_*K$ have the same
cohomology sheaves. We omit the verification that the
canonical base change map $\pi_Y^{-1}Rf_*K \to Rf_{Zar, *}\pi_X^{-1}K$
induces this isomorphism.
\end{proof}

\end{document}
